% appA_lean.tex is GENERATED from this template by gen_appA_r4.py.
% The prose below is hand-edited. Each %%EXTRACT line is replaced by the verbatim lines of the named file of the release tree
% (branch release-v1, read with `git show`, so the working copy cannot interfere) between the first line
% containing the start text and the first line at or after it containing the end text. A path `replay:<file>` means
% ZetaSReplay/<file> of the release tree (the generator records the source of each quoted file in gen_appA_r4_checks.txt).
% A line that begins with %%IF <flag>%% (or %%IFNOT <flag>%%) is kept, without the marker, only if the flag holds (does not hold)
% in the release tree; the generator computes the flags from the quoted files and records them in gen_appA_r4_checks.txt.
% Edit THIS file, then run: python gen_appA_r4.py
\section{The Lean statements}\label{app:lean}
\sloppy

The Lean text below is copied verbatim, by a script that reads the files, from the release of the formalisation: the repository
\repourl{} at tag \repotag{} (\S\ref{app:layout}). Only line breaks inside long lines are added by the typesetting.

\emph{What is claimed} is stated in three files, which import nothing but Mathlib and one another: \texttt{comparator/ChallengeDeps.lean} of the Lean artifact
of~\cite{AF}, which defines the zeros of \texttt{riemannZeta} and the counting functions (\S\ref{app:counting}); the definition layer of this
paper, \texttt{comparator/ChallengeDeps/FollowUpZeta.lean}, with the windows, the local inequalities and the three certificate
\texttt{Prop}s (\S\ref{app:defs}); and the fourteen statements of record, \texttt{comparator/Challenge/FollowUpZeta.lean}
(\S\ref{app:record}). A reader who wants to know what is claimed needs only these three files. The statements of record carry
\texttt{sorry} by design: they are the challenge side of the comparator set-up of the repository of~\cite{AF}. The file
\texttt{comparator/Solution/FollowUpZeta.lean} proves exactly these fourteen statements from the library \texttt{ZetaS}, whose definitions
are copies of those of \S\ref{app:defs} in the namespace \texttt{ZetaS}; the conversions between the two are proofs, not trusted text. The
library's theorems are in \S\ref{app:headline}, the chain for the $K=5$ certificate in \S\ref{app:chain}, the replay's theorems in
\S\ref{app:replay} and the majorant in \S\ref{app:majorant}.

The doc-strings are part of the Lean text and are reproduced unchanged. They are development notes, not part of what is claimed, and
not all current: they refer to an earlier draft of this paper by its labels (for instance \texttt{thm:zeta-mainw2}
is Theorem~\ref{thm:S}, \texttt{thm:sigd-Sigma} and \texttt{thm:sigd-D} are Theorem~\ref{thm:main}, \texttt{cor:oll-SC} is
Corollary~\ref{cor:SC}, \texttt{lem:sigd-env} is Lemma~\ref{lem:room}) and to its line numbers, they name work packages, folders and
dates of the development, and ``level A'' in them means kernel-checked in the sense of \S\ref{ssec:formal}.
%%IF docstring_inprogress%% The doc-string of \texttt{CertAM5} in \S\ref{app:defs} speaks of a kernel replay ``in progress''; that replay is complete (\S\ref{ssec:formal}(c), \S\ref{ssec:replay}, \S\ref{app:replay}).
The data files that the doc-strings name are in \texttt{supplementary/certificates/}, with the same sha256 (Table~\ref{tab:certs}). The
header comment of \S4 of the definition layer says that the certificates were established in Arb ball arithmetic; as \S\ref{ssec:bnb} states,
the box bounds were computed in outward-rounded float64 intervals, with the kernel values enclosed in Arb. Only the Lean code is claimed.

\subsection{The counting functions}\label{app:counting}
From \texttt{comparator/ChallengeDeps.lean} of the repository of~\cite{AF}, unchanged from the tree as shipped (the functions used below are
\texttt{Ncount}, \texttt{Ndist}, \texttt{N0simple}, \texttt{Nsimple} and \texttt{zerosIn}). The phrase ``Not used by the challenge
statements'' in two of these doc-strings refers to the challenge statements of~\cite{AF}; \texttt{Nsimple} is used by those of
\S\ref{app:record}.

%%EXTRACT|comparator/ChallengeDeps.lean|/-! ## 1. Nontrivial zeros|def Nsimple (T₁ T₂ : ℝ)%%

\subsection{The definitions of this paper}\label{app:defs}
The definition layer \texttt{comparator/ChallengeDeps/FollowUpZeta.lean} (namespace \texttt{FollowUpZeta}), from its first line of code to
its end; the file's header comment, which lists its contents, is omitted.

%%EXTRACT|comparator/ChallengeDeps/FollowUpZeta.lean|import ChallengeDeps|end FollowUpZeta%%

\subsection{The fourteen statements of record}\label{app:record}
The statement file \texttt{comparator/Challenge/FollowUpZeta.lean}, from its first line of code to its end (the header comment is omitted).
The six headline statements are those of Theorems~\ref{thm:main} and~\ref{thm:S} and Corollary~\ref{cor:SC}, each with its single
certificate hypothesis; the six \texttt{\_cumulative} forms are the same bounds on the windows $0<\Im\rho\le T$; the last two statements
show that the counts are genuine (every window holds finitely many zeros) and that the denominator tends to infinity.

%%EXTRACT|comparator/Challenge/FollowUpZeta.lean|import ChallengeDeps.FollowUpZeta|end%%

\subsection{The library's headline theorems}\label{app:headline}
The four certificate-only theorems for simple and distinct zeros (\texttt{ZetaS/Top/TopFinal.lean}). They apply the forms with the majorant
threshold $\frac{33}{10}$ to the proved majorant certificate \texttt{MajSound.certMajV2\_proved} (\S\ref{app:majorant}), so that their only
hypothesis is the local certificate:

%%EXTRACT|ZetaS/Top/TopFinal.lean|namespace ZetaS|end ZetaS%%

The two theorems for simple zeros on the line and for zeros that are simple or on the line (\texttt{ZetaS/Top/TopSSC.lean}, namespace \texttt{ZetaS.Top});
statements only, each followed in the file by a tactic proof of about ten lines:

%%EXTRACT|ZetaS/Top/TopSSC.lean|/-- **S — simple zeros on the critical line**|N0simple T (2 * T) := by%%

%%EXTRACT|ZetaS/Top/TopSSC.lean|/-- **SC — simple or critical zeros**|Nsc T (2 * T) := by%%

The statements of these six theorems are, character for character, those of the six dyadic statements of record in \S\ref{app:record},
with the definitions of the namespace \texttt{ZetaS} in place of those of \texttt{FollowUpZeta}; \texttt{comparator/Solution/FollowUpZeta.lean}
derives the fourteen statements of record from them. The library's definition of \texttt{CertAM5} (\texttt{ZetaS/Top/TopDefs.lean}) is:

%%EXTRACT|ZetaS/Top/TopDefs.lean|/-- **Certificate AM5**|LocalCertAM (kPsi psiCos16) W%%

\subsection{The chain for the \texorpdfstring{$K=5$}{K=5} certificate}\label{app:chain}
From \texttt{ZetaS/Cert/ChainV3.lean}: the soundness theorem of the corrected checker and the theorem that twenty checked certificates, one
per canonical class, give \texttt{CertAM5} (the file ends with a diagnostic \texttt{\#eval}, omitted):

%%EXTRACT|ZetaS/Cert/ChainV3.lean|namespace ZetaS.CertV2.ChainV3|end ZetaS.CertV2.ChainV3%%

\subsection{The replay's theorems: the \texorpdfstring{$K=5$}{K=5} pair without hypothesis}\label{app:replay}
These are in the library \texttt{ZetaSReplay} of the repository, which holds the $1{,}478$ files of the replay (\S\ref{ssec:replay}) in one
flat folder, byte-identical to the files that the kernel checked, and is not among the default build targets; \texttt{lake build
ZetaSReplay} repeats the replay (\S\ref{ssec:reproduce}). The replay compiled all $1{,}478$ files with no failure, in $2$~h~$50$~min of
wall time and $20.4$ Lean-hours, against the library at commit
\texttt{bb7223b}, whose Lean files the release changes only in comments (\S\ref{app:layout}). From \texttt{ZetaSReplay/ReplayAM5.lean}, the
statement of \texttt{certAM5\_replayed} and the first line of its proof (the proof term continues with the twenty theorems
\texttt{R<class>.cert\_ok} proved by the kernel; omitted):

%%EXTRACT|replay:ReplayAM5.lean|namespace ZetaS.CertV2|R22222.cert] rfl%%

From \texttt{ZetaSReplay/ReplayHeadlines.lean}, the two theorems of the $K=5$ pair without hypothesis; their statements are those of
\texttt{ZetaS.Top.zeta\_simple\_K5\_final} and \texttt{zeta\_distinct\_K5\_final} (\S\ref{app:headline}) without the hypothesis
\texttt{hcert}:

%%EXTRACT|replay:ReplayHeadlines.lean|namespace ZetaS|end ZetaS%%

\subsection{The majorant}\label{app:majorant}
The majorant is not a hypothesis of any headline theorem: it is proved. Its \texttt{Prop} (\texttt{ZetaS/ChallengeZetaS.lean}, namespace
\texttt{ZetaS}), the instance used (\texttt{ZetaS/Majorant/CertMajV2.lean}) and the theorem (\texttt{ZetaS/Majorant/CertMajV2Proof.lean},
namespace \texttt{ZetaS.MajSound}):

%%EXTRACT|ZetaS/ChallengeZetaS.lean|/-- **`MajorantCert ψ λ`**|2 * (∫ s, B s) + 4 * βstar < lam%%

%%EXTRACT|ZetaS/Majorant/CertMajV2.lean|/-- **CertMajV2**|def CertMajV2 : Prop%%

%%EXTRACT|ZetaS/Majorant/CertMajV2Proof.lean|/-- **CertMajV2**, instantiated|theorem certMajV2_proved%%

\texttt{ZetaS/Majorant/CertMajProof.lean} also proves, as \texttt{MajSound.certMaj\_proved}, the first version \texttt{CertMaj} of
\texttt{ZetaS/ChallengeZetaS.lean}, with the threshold $3.2063638853$; no theorem uses it.
%%IF certmaj_le%% (Its doc-string says $\le$ where the \texttt{Prop} has $<$.)

\subsection{Axiom census and \texttt{sorry}}\label{app:census}
At commit \texttt{bb7223b}, \texttt{\#print axioms} gives exactly \texttt{[propext, Classical.choice, Quot.sound]} for
each of the six theorems of \S\ref{app:headline}, for \texttt{ZetaS.CertV2.ChainV3.certAM5\_of\_checks} and
\texttt{ZetaS.CertV2.ChainV3.cert\_check\_sound}, and for all fourteen theorems of
\texttt{comparator/Solution/FollowUpZeta.lean}. For the replay's three theorems of \S\ref{app:replay} it gives the same three axioms, and
no open declaration lies beneath them. In the repository,
\texttt{comparator/PrintAxioms/FollowUpZeta.lean} prints the axioms of the fourteen theorems of record and
\texttt{comparator/PrintAxioms/FollowUpReplay.lean} those of the three theorems of \S\ref{app:replay} (\S\ref{ssec:reproduce}); the logs and
the \texttt{\#print axioms} outputs of the builds of the release are in \texttt{audit/followup/}.

The library \texttt{ZetaS} has $206$ files and $22{,}751$ lines. It contains $3$ \texttt{sorry}, all in the open node
\texttt{ZetaS/Skeleton/A8g\_SigmaDistGeneral.lean}: Theorems~\ref{thm:OLL} and~\ref{thm:SigmaD} for arbitrary certificate data. That node
is used only by the forms with a general majorant threshold, \texttt{ZetaS.Top.zeta\_simple\_K5\_of\_majorant} and its three companions
(\texttt{ZetaS/Top/TopSolutionGeneral.lean}), and by no headline theorem; the import closure of the fourteen theorems of record contains
no \texttt{sorry}. There is no \texttt{axiom}, \texttt{admit} or \texttt{native\_decide} in the library.

\subsection{Repository layout}\label{app:layout}
The formalisation is published in a modified copy of the Lean repository of~\cite{AF} (\texttt{zeta-23-lean}): \repourl, at tag
\repotag. The commit of the tag is not printed in this paper; it is recorded in the file \texttt{RELEASE\_NOTES.md} of the repository and in
the Zenodo record of the release. The repository is a snapshot of the release without its development history: the commits named in this paper
(such as \texttt{bb7223b}) are commits of that history. The parts that concern this paper:
\begin{itemize}
\item \texttt{ZetaS/} with its root module \texttt{ZetaS.lean}: the library of this paper, outside the default build targets (its folders are
listed below); \texttt{ZetaSCertShims/}: three forwarding modules through which the checker files are imported by their bare names, so
that they are kept byte for byte as checked;
\item \texttt{ZetaSReplay/}: the $1{,}478$ files of the kernel replay of \texttt{CertAM5} ($1{,}475$ generated shards and the three common
modules \texttt{ReplayCommon}, \texttt{ReplayAM5} and \texttt{ReplayHeadlines}), in one flat folder, as a library outside the default build
targets (\S\ref{ssec:replay}, \S\ref{app:replay});
\item \texttt{comparator/}: the comparator set-up of~\cite{AF}, with seven new files for this paper: \texttt{ChallengeDeps/FollowUpZeta.lean} and
\texttt{Challenge/FollowUpZeta.lean} (quoted above), \texttt{Solution/FollowUpZeta.lean}, \texttt{PrintAxioms/FollowUpZeta.lean} and
\texttt{PrintAxioms/FollowUpReplay.lean}, \texttt{config-followup.json} and \texttt{README\_followup.md};
\item \texttt{supplementary/}: the data, certifiers and logs of Table~\ref{tab:certs} and the generator of the replay files
(\S\ref{ssec:reproduce});
\item \texttt{papers/zeta/}: the sources and the PDF of this paper; \texttt{audit/followup/}: the logs and \texttt{\#print axioms} outputs of
the builds of the release.
\end{itemize}
The repository also contains the library \texttt{ZetaShell/} of a companion paper on families of Dirichlet $L$-functions, whose sources are in
\texttt{papers/family/}, and the library \texttt{ZetaQ/} added to the repository earlier (see \texttt{NOTICE}); \texttt{ZetaS} imports neither.
The statement files of~\cite{AF} (\texttt{Zeta23/Statement.lean} and the existing files of \texttt{comparator/}) are unchanged. Further,
the repository contains changes to $25$ files of the library \texttt{Zeta23} ($3{,}122$ lines added, $121$ deleted, between the upstream base
commit and \texttt{bb7223b}): a weaker parameter condition
(\texttt{Params.ValidQ} in place of \texttt{Params.Valid}), with the proofs adapted, and new modules for Gevrey ramps and a free buffer
(\texttt{Tail/GevreyTail}, \texttt{Taper/GevreyProduct}, \texttt{WindowD}). The axiom census above is for the tree with these changes.
The Lean files of \texttt{Zeta23}, \texttt{ZetaS}, \texttt{ZetaSCertShims} and \texttt{comparator/} in the release are those of commit
\texttt{bb7223b}, at which the census of \S\ref{app:census} was taken and against which the replay ran, except for comments: for the release,
the copyright headers of the four new comparator files of commit \texttt{bb7223b} and stale doc-strings in
\texttt{comparator/ChallengeDeps/FollowUpZeta.lean} and \texttt{ZetaS/ChallengeZetaS.lean} were corrected. Two files were added to
\texttt{comparator/PrintAxioms/}, \texttt{FollowUpReplay.lean} (\S\ref{app:census}) and \texttt{FollowUpShell.lean}, for the companion paper,
and \texttt{comparator/README\_followup.md} was brought up to date. The comparator package \texttt{FollowUpFamily} of the companion paper
(\texttt{ChallengeDeps/}, \texttt{Challenge/}, \texttt{Solution/} and \texttt{PrintAxioms/FollowUpFamily.lean}, \texttt{config-followup-family.json}
and \texttt{README\_followup\_family.md}) was also added.
The folders of \texttt{ZetaS}:
\begin{itemize}
\item \texttt{ZetaS/Interfaces.lean}, \texttt{InterfacesV2.lean}: the interfaces between the zero side and the local arguments (Gram data,
frame families at bandwidth $\lambda<1$); \texttt{ZetaS/ChallengeZetaS.lean}: the library's copies of the definitions of \S\ref{app:defs},
together with \texttt{MajorantCert} and \texttt{CertMaj};
\item \texttt{LinAlg/}: Section~\ref{sec:linalg} (stability inequality and its two-parameter form, $\Phi_m$, pinching, monotonicity, pair
removal, localisation, inertia);
\item \texttt{Window/}: the two windows, their exact moments and constants, and the window-general pipeline at bandwidth $\lambda<1$;
\texttt{Bandwidth/}: the limit $\lambda\to1^-$;
\item \texttt{ZeroSide/}: the Poisson and Gram identities, positive definiteness, truncation defect, inertia of pair blocks, and the assembly
from the zeros of \texttt{riemannZeta};
\item \texttt{SigmaDist/}: Section~\ref{sec:simple}; \texttt{KSide/}: Sections~\ref{sec:first} and~\ref{sec:sc}; \texttt{Top/}: the numerics of
the constants, the robustness conditions, the headline theorems (\texttt{TopFinal.lean}, \texttt{TopSSC.lean}) and the forms with a general
majorant threshold (\texttt{TopSolutionGeneral.lean});
\item \texttt{Majorant/}: the majorant, its analytic part, its cell checker, the $13$ replay shards and the two majorant theorems;
\texttt{Cert/}: the corrected checker for the $K=5$ certificate (\texttt{CheckerBase.lean}, \texttt{CheckerLeaves.lean},
\texttt{CheckerCoreV3.lean}, with kernel-checked negative controls in \texttt{CheckerTestsV3.lean}), the specification
\texttt{CertSpecAM5.lean}, the soundness nodes and the chain \texttt{ChainV3.lean};
\item \texttt{Skeleton/}: re-export files for the proved nodes, and the one open node \texttt{A8g\_SigmaDistGeneral.lean}.
\end{itemize}
The new Lean files of \texttt{comparator/} carry the author's copyright notice under the Apache License~2.0, the licence of the repository
of~\cite{AF}; the files that come from~\cite{AF} keep their own notices (see \texttt{NOTICE}).
